% GENERATED FILE. Edit canonical lessons, then run npm run generate:books.
% canonical-source-hash: fnv1a64:662c1032c0f8455a
% canonical-lessons: AR-C501-hakim, AR-C501-mahir, AR-C501-mubdi, AR-C501-aqil, AR-C501-thiqa

\chapter{Wise, skilled and sure}
\label{ch:ar-wise-skilled-and-sure}
\hlchaptermodality{501}

\begin{chapteropening}
\textbf{By the end of this chapter:} \emph{I can describe a person as wise, skilful, creative or sensible, and talk about confidence.}

Complete the last lesson of chapter 501: I can describe a person as wise, skilful, creative or sensible, and talk about confidence.
\end{chapteropening}

\section[\texorpdfstring{\ar{حكيم} (ḥakīm)}{حكيم (ḥakīm)}]{\ar{حكيم} (ḥakīm) — wise}

\label{lesson:AR-C501-hakim}



\begin{quote}
Before the new one: say the Arabic for cowardly, then the Arabic for bold.
\end{quote}

\subsection*{You'll want to know: \ar{حكيم}}
\textbf{\ar{حكيم}} — \emph{ḥakīm} — \textquotedblleft{}wise\textquotedblright{}.

\begin{quote}
\ar{جدتي} \ar{امرأة} \ar{حكيمة}.
\end{quote}

\begin{quote}
\emph{jaddatī imraʾa ḥakīma.}
\end{quote}

\begin{quote}
My grandmother is a wise woman.
\end{quote}

In some spoken dialects \ar{حكيم} is also a doctor (colloquial).

\begin{cousinweb}
\textbf{\ar{حكيم}} (\emph{ḥakīm}) is built on the native Arabic root \emph{ḥ-k-m}, \textquotedblleft{}to judge, to restrain\textquotedblright{} (Lane). The noun ḥikma is wisdom.
\end{cousinweb}

\begin{grammarlens}[title={What you've built}]
One more word for describing what you lived and felt.
\end{grammarlens}

\subsection*{Your turn}
\begin{itemize}
\raggedright
  \item \emph{Say it:} \emph{ḥakīm}
  \item \emph{Say it:} \emph{ḥakīm}, once more
  \item \emph{Say it:} \emph{jaddatī imraʾa ḥakīma.} — My grandmother is a wise woman.
  \item \emph{Recall:} say the Arabic for cowardly, then the Arabic for bold, then say \emph{ḥakīm} again
\end{itemize}

\begin{tcolorbox}[breakable,colback=teal!4,colframe=teal!35!black,title={Before you move on}]
What does \ar{حكيم} mean? (\textquotedblleft{}Wise\textquotedblright{}.) Say it once more.
\end{tcolorbox}
\section[\texorpdfstring{\ar{ماهر} (māhir)}{ماهر (māhir)}]{\ar{ماهر} (māhir) — skilful, skilled}

\label{lesson:AR-C501-mahir}



\begin{quote}
Before the new one: say the Arabic for bold, then the Arabic for wise.
\end{quote}

\subsection*{You'll want to know: \ar{ماهر}}
\textbf{\ar{ماهر}} — \emph{māhir} — \textquotedblleft{}skilful, skilled\textquotedblright{}.

\begin{quote}
\ar{صنع} \ar{النجار} \ar{الماهر} \ar{طاولة} \ar{جميلة}.
\end{quote}

\begin{quote}
\emph{ṣanaʿa an-najjār al-māhir ṭāwila jamīla.}
\end{quote}

\begin{quote}
The skilful carpenter made a beautiful table.
\end{quote}

The noun \ar{مهارة}, a skill, shares the root.

\begin{cousinweb}
\textbf{\ar{ماهر}} (\emph{māhir}) is built on the native Arabic root \emph{m-h-r}, \textquotedblleft{}to be skilled\textquotedblright{} (Lane). Māhir is the active participle.
\end{cousinweb}

\begin{grammarlens}[title={What you've built}]
One more word for describing what you lived and felt.
\end{grammarlens}

\subsection*{Your turn}
\begin{itemize}
\raggedright
  \item \emph{Say it:} \emph{māhir}
  \item \emph{Say it:} \emph{māhir}, once more
  \item \emph{Say it:} \emph{ṣanaʿa an-najjār al-māhir ṭāwila jamīla.} — The skilful carpenter made a beautiful table.
  \item \emph{Recall:} say the Arabic for bold, then the Arabic for wise, then say \emph{māhir} again
\end{itemize}

\begin{tcolorbox}[breakable,colback=teal!4,colframe=teal!35!black,title={Before you move on}]
What does \ar{ماهر} mean? (\textquotedblleft{}Skilful, skilled\textquotedblright{}.) Say it once more.
\end{tcolorbox}
\section[\texorpdfstring{\ar{مبدع} (mubdiʿ)}{مبدع (mubdiʿ)}]{\ar{مبدع} (mubdiʿ) — creative}

\label{lesson:AR-C501-mubdi}



\begin{quote}
Before the new one: say the Arabic for wise, then the Arabic for skilful.
\end{quote}

\subsection*{You'll want to know: \ar{مبدع}}
\textbf{\ar{مبدع}} — \emph{mubdiʿ} — \textquotedblleft{}creative\textquotedblright{}.

\begin{quote}
\ar{ابنتي} \ar{طالبة} \ar{مبدعة} \ar{في} \ar{الرسم}.
\end{quote}

\begin{quote}
\emph{ibnatī ṭāliba mubdiʿa fī ar-rasm.}
\end{quote}

\begin{quote}
My daughter is a creative student in drawing.
\end{quote}

\begin{cousinweb}
\textbf{\ar{مبدع}} (\emph{mubdiʿ}) is built on the native Arabic root \emph{b-d-ʿ}, \textquotedblleft{}to originate, to make a thing new\textquotedblright{} (Lane). Mubdiʿ is the participle of Form IV: one who makes a new thing.
\end{cousinweb}

\begin{grammarlens}[title={What you've built}]
One more word for describing what you lived and felt.
\end{grammarlens}

\subsection*{Your turn}
\begin{itemize}
\raggedright
  \item \emph{Say it:} \emph{mubdiʿ}
  \item \emph{Say it:} \emph{mubdiʿ}, once more
  \item \emph{Say it:} \emph{ibnatī ṭāliba mubdiʿa fī ar-rasm.} — My daughter is a creative student in drawing.
  \item \emph{Recall:} say the Arabic for wise, then the Arabic for skilful, then say \emph{mubdiʿ} again
\end{itemize}

\begin{tcolorbox}[breakable,colback=teal!4,colframe=teal!35!black,title={Before you move on}]
What does \ar{مبدع} mean? (\textquotedblleft{}Creative\textquotedblright{}.) Say it once more.
\end{tcolorbox}
\section[\texorpdfstring{\ar{عاقل} (ʿāqil)}{عاقل (ʿāqil)}]{\ar{عاقل} (ʿāqil) — sensible}

\label{lesson:AR-C501-aqil}



\begin{quote}
Before the new one: say the Arabic for skilful, then the Arabic for creative.
\end{quote}

\subsection*{You'll want to know: \ar{عاقل}}
\textbf{\ar{عاقل}} — \emph{ʿāqil} — \textquotedblleft{}sensible\textquotedblright{}.

\begin{quote}
\ar{تكلم} \ar{مثل} \ar{رجل} \ar{عاقل}.
\end{quote}

\begin{quote}
\emph{takallama mithla rajul ʿāqil.}
\end{quote}

\begin{quote}
He spoke like a sensible man.
\end{quote}

The noun \ar{عقل} is the mind.

\begin{cousinweb}
\textbf{\ar{عاقل}} (\emph{ʿāqil}) is built on the native Arabic root \emph{ʿ-q-l}, \textquotedblleft{}to bind; to understand\textquotedblright{} (Lane). Maʿqūl, reasonable, which you know, is from the same root.
\end{cousinweb}

\begin{grammarlens}[title={What you've built}]
One more word for describing what you lived and felt.
\end{grammarlens}

\subsection*{Your turn}
\begin{itemize}
\raggedright
  \item \emph{Say it:} \emph{ʿāqil}
  \item \emph{Say it:} \emph{ʿāqil}, once more
  \item \emph{Say it:} \emph{takallama mithla rajul ʿāqil.} — He spoke like a sensible man.
  \item \emph{Recall:} say the Arabic for skilful, then the Arabic for creative, then say \emph{ʿāqil} again
\end{itemize}

\begin{tcolorbox}[breakable,colback=teal!4,colframe=teal!35!black,title={Before you move on}]
What does \ar{عاقل} mean? (\textquotedblleft{}Sensible\textquotedblright{}.) Say it once more.
\end{tcolorbox}
\section[\texorpdfstring{\ar{ثقة} (thiqa)}{ثقة (thiqa)}]{\ar{ثقة} (thiqa) — confidence; trust}

\label{lesson:AR-C501-thiqa}



\begin{quote}
Before the new one: say the Arabic for creative, then the Arabic for sensible.
\end{quote}

\subsection*{You'll want to know: \ar{ثقة}}
\textbf{\ar{ثقة}} — \emph{thiqa} — \textquotedblleft{}confidence; trust\textquotedblright{}.

\begin{quote}
\ar{تتكلم} \ar{أختي} \ar{الإنجليزية} \ar{بثقة}.
\end{quote}

\begin{quote}
\emph{tatakallamu ukhtī al-injilīziyya bi-thiqa.}
\end{quote}

\begin{quote}
My sister speaks English with confidence.
\end{quote}

\ar{ثقة} \ar{بالنفس} is self-confidence.

\begin{cousinweb}
\textbf{\ar{ثقة}} (\emph{thiqa}) is built on the native Arabic root \emph{w-th-q}, \textquotedblleft{}to trust, to rely on\textquotedblright{} (Lane). The verb wathiqa, to trust, which you know, is from the same root; thiqa drops the first letter w.
\end{cousinweb}

\begin{grammarlens}[title={What you've built}]
One more word for describing what you lived and felt.
\end{grammarlens}

\subsection*{Your turn}
\begin{itemize}
\raggedright
  \item \emph{Say it:} \emph{thiqa}
  \item \emph{Say it:} \emph{thiqa}, once more
  \item \emph{Say it:} \emph{tatakallamu ukhtī al-injilīziyya bi-thiqa.} — My sister speaks English with confidence.
  \item \emph{Recall:} say the Arabic for creative, then the Arabic for sensible, then say \emph{thiqa} again
\end{itemize}

\begin{tcolorbox}[breakable,colback=teal!4,colframe=teal!35!black,title={Before you move on}]
What does \ar{ثقة} mean? (\textquotedblleft{}Confidence; trust\textquotedblright{}.) Say it once more.
\end{tcolorbox}
